\section{Introduction}

Single-cell RNA sequencing (scRNA-seq) has made it possible to resolve heterogeneous tissues at cellular scale and to construct reference atlases spanning organs, developmental stages and disease states \citep{Macosko2015,Zheng2017,Tanay2017,Regev2017,TabulaSapiens2022,LueckenTheis2019}. Interpreting these data still depends on cell-type annotation: the assignment of biologically meaningful identities to measured expression profiles. Manual annotation based on canonical markers remains valuable, but it is laborious, sensitive to expert judgement and difficult to scale. Automated methods therefore occupy a central place in atlas construction and reference mapping \citep{Aran2019,Kiselev2018,AlquicirayHernandez2019,Hao2021,Abdelaal2019,Pasquini2021}.

The available annotation strategies span a wide range of statistical complexity. Logistic regression, support-vector machines, nearest-neighbour projection, tree ensembles and reference-mapping systems can perform strongly when labelled cells from an appropriate reference population are available \citep{Aran2019,Kiselev2018,AlquicirayHernandez2019,DominguezConde2022}. Latent-variable models such as scVI and scANVI additionally provide probabilistic representations and batch-aware integration \citep{Lopez2018,Luecken2022}. These methods establish demanding baselines: a more complex representation is useful for supervised annotation only if it preserves task-relevant biological signal, transfers across the intended biological partition and improves an endpoint that matters in practice.

Single-cell foundation models seek to learn reusable representations from very large transcriptomic corpora. scBERT, Geneformer, scGPT, scFoundation and Universal Cell Embeddings differ in tokenization, architecture and pretraining objective, but share the premise that large-scale pretraining can encode regularities that are useful across downstream tasks \citep{Yang2022,Theodoris2023,Cui2024,Hao2024,Rosen2023}. scGPT is a generative transformer pretrained on more than 33 million cells and has been applied to cell-type annotation, integration, perturbation response and multi-omic analysis \citep{Cui2024}. Its released whole-human checkpoint and public embedding interface make it a reproducible case study for asking whether a fixed pretrained cell representation improves supervised annotation.

The answer depends on the evaluation regime. Frozen embeddings followed by a supervised classifier are neither zero-shot annotation nor end-to-end fine-tuning. They test whether a fixed representation makes target labels accessible to a downstream learner. Strong performance in this regime may be valuable, but it does not establish transfer to unlabelled cell types, external atlases or biological conditions absent from the target training data. Conversely, weak frozen-embedding performance does not rule out benefits from official fine-tuning, other checkpoints or other foundation models. Recent independent evaluations have similarly emphasized that foundation-model conclusions depend on task definition, pretraining exposure, comparator choice and adaptation protocol \citep{Boiarsky2024,Kedzierska2025}.

Biological partitioning is especially important in single-cell benchmarks. A random cell split allows cells from the same donor, sample or processing batch to appear in both training and testing. Such a split measures interpolation among cells from an atlas, but can yield optimistic estimates of deployment to unseen biological units. Grouped validation instead keeps all cells from a donor together, preserving the hierarchy of the data and making the test question explicit \citep{Roberts2017,Squair2021}. It does not by itself create external-atlas validation---cohort-specific acquisition, curation and label definitions remain shared---but it removes direct donor overlap and provides multiple biological test partitions.

Comparator construction also changes the scientific question. scGPT returns a 512-dimensional, unit-normalized cell embedding, whereas a direct expression baseline can contain more than one thousand sparse, log-normalized gene values on a different numerical scale. Applying the same regularization setting to both representations need not be neutral. Hyperparameters should therefore be selected separately inside the training data for every representation--classifier pair. A dimensionality-matched expression control is also informative: if direct expression outperforms a 512-dimensional embedding, a 512-component expression projection can help distinguish a scGPT-specific effect from the generic information loss associated with compression.

Annotation systems are often used through their confidence scores as well as their discrete predictions. Accuracy and Macro F1 assess discrimination, whereas calibration asks whether predictive probabilities agree with observed frequencies \citep{Guo2017,Vaicenavicius2019,SilvaFilho2023}. These properties are distinct. Expected calibration error (ECE) summarizes a reliability diagram but depends on binning; the multiclass Brier score and negative log-likelihood are proper scoring rules that jointly reflect probability quality and discrimination \citep{Brier1950,Gneiting2007,Nixon2019}. Moreover, probabilities in a frozen-embedding pipeline are generated by the downstream classifier. Calibration findings must consequently be attached to complete representation--classifier pipelines, not attributed to the encoder alone. Post-hoc temperature scaling provides a simple test of how much probability error can be corrected without changing the predicted class \citep{Guo2017}.

We conducted an accuracy and calibration benchmark of official frozen scGPT embeddings, matched normalized expression and a 512-component expression projection. Frozen scGPT and matched expression were each paired with Logistic Regression, Random Forest and XGBoost, with hyperparameters selected independently inside the training donors. The primary evaluation used five donor-held-out folds in three human atlases spanning breast cancer, peripheral blood and brain. SVD-512 was evaluated as a dimensionality control with Logistic Regression in all three atlases and with all three classifiers in the TNBC and brain atlases. We further examined class-specific errors, labelled-cell scarcity with labelled-only feature selection, held-out-donor temperature scaling in two fully audited atlases and a separately reported pig-pancreas stress test using the human checkpoint. This design addresses a specific question: under supervised annotation of broad portal labels, how do fixed scGPT 0.2.5 embeddings compare with equally trained classical representations when the evaluated cells come from donors not used to fit the classifier?
